% -*- coding: utf-8 -*-
% !TEX program = xelatex
\documentclass[plain,noanswer,medmath]{../../template/jnuexam}
\usepackage{../../template/kysx}

\begin{document}


\examtitle{1992}{1}


\makepart{填空题}{本题共5小题，每小题3分，满分15分}


\begin{problem}
设函数$y = y(x)$由方程$\e^{x + y} + \cos(xy) = 0$确定，则$\frac{\dy}{\dx} =$ \fillin{}．
\end{problem}


\begin{problem}
函数$u = \ln(x^2 + y^2 + z^2)$在点$M(1,2, -2)$处的梯度$\left.\grad u\right|_M =$ \fillin{}．
\end{problem}


\begin{problem}
设$f(x) = \begin{cases}
      -1, & - \pi < x \le 0, \\
 1 + x^2, & 0 < x \le \pi , \\
\end{cases}$则其以$2\pi$为周期的傅里叶级数在点$x = \pi$处收敛于 \fillin{}．
\end{problem}


\begin{problem}
微分方程$y' + y\tan x = \cos x$的通解为$y =$ \fillin{}．
\end{problem}


\begin{problem}
设$A = \begin{pmatrix}
  a_1b_1 & a_1b_2 & \cdots & a_1b_n \\
  a_2b_1 & a_2b_2 & \cdots & a_2b_n \\
 \vdots & \vdots &        & \vdots \\
  a_nb_1 & a_nb_2 & \cdots & a_nb_n \\
\end{pmatrix}$，其中$a_i \ne 0$,$b_i \ne 0$($i = 1,2 \cdots n$)．则矩阵$A$的秩$r(A)=$ \fillin{}．
\end{problem}


\makepart{选择题}{本题共5小题，每小题3分，满分15分}


\begin{problem}
当$x \to 1$时，函数$\frac{x^2 - 1}{x - 1}\e^{\frac{1}{x - 1}}$的极限 \pickout{}
\begin{abcd}
\item 等于$2$．     
\item 等于$0$．     
\item 为$\infty$．     
\item 不存在但不为$\infty$．
\end{abcd}
\end{problem}


\begin{problem}
级数$\sum_{n = 1}^{\infty}(- 1)^n\left(1 - \cos \frac{\alpha}{n} \right)$(常数$\alpha > 0$) \pickout{}
\begin{abcd}
\item 发散．   
\item 条件收敛．   
\item 绝对收敛．   
\item 收敛性与$\alpha$有关．
\end{abcd}
\end{problem}


\begin{problem}
在曲线$x = t,y = - t^2,z = t^3$的所有切线中，与平面$x + 2y + z = 4$平行的切线 \pickout{}
\begin{abcd}
\item 只有1条．    
\item 只有2条．    
\item 至少有3条．    
\item 不存在．
\end{abcd}
\end{problem}


\begin{problem}
设$f(x) = 3x^3 + x^2|x|$，则使$f^n(0)$存在的最高阶数$n$为 \pickout{}
\begin{abcd}
\item $0$．          
\item $1$．          
\item $2$．          
\item $3$．
\end{abcd}
\end{problem}


\begin{problem}
要使$\xi_1=\begin{pmatrix}1 \\ 0 \\ 2\end{pmatrix}$, $\xi_2=\begin{pmatrix}0 \\ 1 \\ -1\end{pmatrix}$
都是线性方程组$Ax=0$的解，只要系数矩阵$A$为\pickout{}
\begin{abcd}
\item $\begin{pmatrix}-2 & 1 &  1\end{pmatrix}$．					
\item $\begin{pmatrix} 2 & 0 & -1 \\ 0 &  1 &  1\end{pmatrix}$．
\item $\begin{pmatrix}-1 & 0 &  2 \\ 0 &  1 & -1\end{pmatrix}$．					
\item $\begin{pmatrix} 0 & 1 & -1 \\ 4 & -2 & -2 \\ 0 & 1 & 1\end{pmatrix}$．
\end{abcd}
\end{problem}


\makepart{计算题}{本题共3小题，每小题5分，满分15分}


\begin{problem}
求$\lim_{x \to 0} \frac{\e^x - \sin x - 1}{1 - \sqrt{1 - x^2}}$．
\end{problem}


\begin{problem}
设$z = f(\e^x\sin y, x^2 + y^2)$，其中$f$具有二阶连续偏导数，
求$\frac{\pd^2 z}{\pdx\pd y}$．
\end{problem}


\begin{problem}
设$f(x) = \begin{cases}
  1 + x^2, & x \le 0, \\
 \e^{- x}, & x > 0, \\
\end{cases}$求$\int_1^3f(x - 2) \dx$．
\end{problem}


\makepart{}{本题满分6分}

\begin{problem}
求微分方程$y'' + 2y' - 3y = \e^{- 3x}$的通解．
\end{problem}


\makepart{}{本题满分8分}

\begin{problem}
计算曲面积分$\iint_{\Sigma}(x^3 + az^2)\dy\dz + (y^3 + ax^2)\dz\dx + (z^3 + ay^2)\dx\dy$，其中$\Sigma$为上半球面$z = \sqrt{a^2 - x^2 - y^2}$的上侧．
\end{problem}


\makepart{}{本题满分7分}

\begin{problem}
设$f''(x) < 0$, $f(0) = 0$，证明对任何$x_1 > 0,x_2 > 0$，有$f(x_1 + x_2) < f(x_1) + f(x_2)$．
\end{problem}


\makepart{}{本题满分8分}

\begin{problem}
在变力$\vec{F}= yz\vec{i}+ zx\vec{j}+ xy\vec{k}$的作用下，质点由原点沿直线运动到椭球面$\frac{x^2}{a^2} + \frac{y^2}{b^2} + \frac{z^2}{c^2} = 1$上第一卦限的点$M(\xi ,\eta ,\zeta)$．问当$\xi ,\eta ,\zeta$取何值时，力$\vec{F}$所做的功$W$最大？并求出$W$的最大值．
\end{problem}


\makepart{}{本题满分7分}

\begin{problem}
设向量组$\alpha_1,\;\alpha_2,\;\alpha_3$线性相关，
向量组$\alpha_2,\;\alpha_3,\;\alpha_4$线性无关，问：
\begin{enumerate}
\item $\alpha_1$能否由$\alpha_2,\;\alpha_3$线性表示？证明你的结论．
\item $\alpha_4$能否由$\alpha_1,\;\alpha_2,\;\alpha_3$线性表示？证明你的结论．
\end{enumerate}
\end{problem}


\makepart{}{本题满分7分}

\begin{problem}
设$3$阶矩阵$A$的特征值为$\lambda_1 = 1,\lambda_2 = 2,\lambda_3 = 3$，对应的特征向量依次为
$$\xi_1 = \begin{pmatrix}
 1 \\
 1 \\
 1 \\
\end{pmatrix},\xi_2 = \begin{pmatrix}
 1 \\
 2 \\
 4 \\
\end{pmatrix},\xi_3 = \begin{pmatrix}
 1 \\
 3 \\
 9 \\
\end{pmatrix},$$
又向量$\beta = \smash[t]{\begin{pmatrix}
 1 \\
 1 \\
 3 \\
\end{pmatrix}}$．
\begin{enumerate*}
\item 将$\beta$用$\xi_1,\xi_2,\xi_3$线性表示．
\item 求$A^n\beta$（$n$为自然数）．
\end{enumerate*}
\end{problem}


\makepart{填空题}{本题共2小题，每小题3分，满分6分}


\begin{problem}
已知$P(A) = P(B) = P(C) = \frac{1}{4}$,$P(AB) = 0$,$P(AC) = P(BC) = \frac{1}{16}$，则事件$A$、$B$、
$C$全不发生的概率为 \fillin{}．
\end{problem}


\begin{problem}
设随机变量$X$服从参数为$1$的指数分布，则数学期望$E(X + \e^{-2X}) =$ \fillin{}．
\end{problem}


\makepart{}{本题满分6分}

\begin{problem}
设随机变量$X$与$Y$独立，$X$服从正态分布$N(\mu ,\sigma^2)$，$Y$服从$[- \pi ,\pi]$上的均匀分布，
试求$Z = X + Y$的概率分布密度（计算结果用标准正态分布函数$\Phi(x)$表示，
其中$\Phi(x) = \frac{1}{\sqrt{2\pi}}\int_{-\infty}^x \e^{- \frac{t^2}{2}} \dt$）．
\end{problem}


\end{document}
